\begin{thebibliography}{99}
\bibitem{FGA221} A.~Grothendieck, \emph{Techniques de construction et th\'eor\`emes d'existence en g\'eom\'etrie alg\'ebrique IV: les sch\'emas de Hilbert}, S\'eminaire Bourbaki, Expos\'e221 (1961), pp.249--276. Lemma2.4, p.255; Theorems3.1--3.2, pp.259--260. \url{https://www.numdam.org/article/SB_1960-1961__6__249_0.pdf}.
\bibitem{EGAIV2} A.~Grothendieck and J.~Dieudonn\'e, \emph{\'El\'ements de g\'eom\'etrie alg\'ebrique IV, deuxi\`eme partie}, Publ. Math. IH\'ES24 (1965), Theorem6.9.1, p.154. \url{https://www.numdam.org/articles/10.1007/BF02684322/}.
\bibitem{SGA4} M.~Artin, A.~Grothendieck and J.-L.~Verdier, \emph{Th\'eorie des topos et cohomologie \'etale des sch\'emas}, SGA4, Expos\'eXVI, Corollary2.2, original p.214, with the locally-constant-coefficient correction in SGA$4\frac12$, p.312. Re-edition: \url{https://www.normalesup.org/~forgogozo/SGA4/16/16.pdf}.
\bibitem{SGA45Cycle} A.~Grothendieck, written by P.~Deligne, \emph{La classe de cohomologie associ\'ee \`a un cycle}, SGA$4\frac12$, LNM569 (1977), pp.129--153, \S\S2.3.1--2.3.3, Lemma2.3.7 and Theorem2.3.8. Retypeset source: \url{https://github.com/NomiL/sga4.5/blob/master/sga4.5-4.tex}.
\bibitem{SGA1} A.~Grothendieck et al., \emph{Rev\^etements \'etales et groupe fondamental}, SGA1, Expos\'eXII, Proposition2.4 and Theorem5.1 (original p.333 for the latter). Re-edition: \url{https://arxiv.org/abs/math/0206203}.
\bibitem{Deligne82} P.~Deligne, notes by J.~S.~Milne, \emph{Hodge cycles on abelian varieties}, in \emph{Hodge Cycles, Motives, and Shimura Varieties}, LNM900 (1982), Proposition2.9(b), Theorem2.11. Corrected author text: \url{https://www.jmilne.org/math/Documents/Deligne82.pdf}.
\bibitem{BrunjesSerpe} L.~Br\"unjes and C.~Serp\'e, \emph{Uniform bounds and ultraproducts of cycles}, arXiv:0901.4853, Definition2.3(ii), Theorem3.4 and the number-field discussion at the end of the introduction. \url{https://arxiv.org/abs/0901.4853}.
\bibitem{OAI032} OpenAI, \emph{The rational Hodge conjecture for CM abelian varieties}, 6 October2026, supplied PDF. Theorem1.1; Corollaries8.3--8.4; Remark8.5, pp.50--51. Treated as a cited input, as instructed by the task.
\bibitem{OAI001} OpenAI, \emph{Milne's rationality conjecture for abelian varieties}, 7 October2026. Equations(3.15)--(3.16), p.16; Corollary3.4, p.17; same-dimensional caution and Corollary6.4, p.33. Official PDF at commit \texttt{fd4aeeb2ee4fc729c18d98444fed42fd0529eeeb}: \url{https://github.com/openai/math/tree/fd4aeeb2ee4fc729c18d98444fed42fd0529eeeb/preprints/Milnes-rationality-conjecture-for-abelian-varieties-October-7-2026}.
\bibitem{MilneReport2026} J.~S.~Milne, \emph{On the Hodge, Tate, and standard conjectures for abelian varieties}, Hexagon2611.00244v1, supplied PDF; \S\S1.7--1.9, \S3.6. The sign in the primitive pairing of \S3.6 is audited separately below; the theorem input is not re-proved here.
\bibitem{VanHoftenSempliner} P.~van Hoften and A.~Sempliner, \emph{The Hodge Conjecture for CM abelian varieties}, Hexagon2610.00259v2, supplied PDF. Carried as the task's companion theorem input; no independent whole-paper verification claimed.
\bibitem{DeligneOrdinary} P.~Deligne, \emph{Vari\'et\'es ab\'eliennes ordinaires sur un corps fini}, Invent. Math.8 (1969), 238--243, \S\S3,7,8. Primary article in translation: \url{https://translations.thosgood.net/IM-8-1969-238.html}.
\bibitem{Tate1966} J.~Tate, \emph{Endomorphisms of abelian varieties over finite fields}, Inventiones Mathematicae2 (1966), 134--144. \S2, pp.138--139; Theorem2(a), p.140. \url{https://pazuki.perso.math.cnrs.fr/index_fichiers/Tate66.pdf}.
\bibitem{MilneAV} J.~S.~Milne, \emph{Abelian Varieties}, corrected2 January2022. Corollary2.2 and Theorem3.1, pp.3--4; Corollaries6.4--6.6 and Theorem6.7; Theorem7.3, p.11; Proposition8.1 and Theorem8.2; \S11, pp.17--18; Theorem13.3, p.25; Proposition16.8 and Remark16.9, p.33. \url{https://www.jmilne.org/math/xnotes/AVs.pdf}.
\bibitem{GriffithsPoincare1982} P.~A.~Griffiths, \emph{Poincar\'e and algebraic geometry}, Bulletin of the American Mathematical Society (New Series)6, no.2 (1982), 147--160. Poincar\'e's formula, p.149, displayed (6)--(7). \url{https://publications.ias.edu/sites/default/files/poincare.pdf}.
\bibitem{MilneEC} J.~S.~Milne, \emph{Elliptic Curves}, second edition, PropositionII.6.4, p.71; TheoremII.6.9, p.73; PropositionIV.7.13, p.140; ChapterV, \S7, p.216. \url{https://www.jmilne.org/math/Books/EC2.pdf}.
\bibitem{Fulton} W.~Fulton, \emph{Intersection Theory}, second edition, Springer1998. Refined pullbacks and projection formula, Chapters6 and8; Example8.4.6, p.146, the general projective B\'ezout degree bound for intersection components. These are classical background results, not new proofs in this run.
\bibitem{Hartshorne} R.~Hartshorne, \emph{Algebraic Geometry}, Springer1977, ChapterV, Theorem1.9, algebraic Hodge index for a smooth projective surface.
\bibitem{SilvermanAEC} J.~H.~Silverman, \emph{The Arithmetic of Elliptic Curves}, second edition, Springer2009. ChapterIII, Theorem10.1 and Corollary10.2; ChapterV, Theorem4.1(a), supersingularity criterion via the Hasse invariant.
\bibitem{CHVW2022} W.~Castryck, J.~Houben, F.~Vercauteren and B.~Wesolowski, \emph{On the decisional Diffie--Hellman problem for class group actions on oriented elliptic curves}, Research in Number Theory8 (2022), article99, footnote2 (characteristic-two Hurwitz order). \url{https://doi.org/10.1007/s40993-022-00399-6}.
\bibitem{Voight} J.~Voight, \emph{Quaternion Algebras}, version1.0.5 (2024), Lemma11.1.2, \S11.2, Theorem11.5.14; Theorem15.5.5, p.241; \S42.1.3 and Theorem42.1.9. \url{https://jvoight.github.io/quat-book-v1.0.5.pdf}.
\bibitem{Schoen1988} C.~Schoen, \emph{Hodge classes on self-products of a variety with an automorphism}, Compositio Math.65 (1988), 3--32. Theorem2.0 and its $r=0$ construction, pp.12--13; \S3, Corollary3.1, Theorem3.2 and actual Prym-dominance calculation, pp.24--30. The missing polarization qualification in Theorem3.2 is corrected by the 1998 addendum. \url{https://www.numdam.org/article/CM_1988__65_1_3_0.pdf}.
\bibitem{Schoen1998} C.~Schoen, \emph{Addendum to: Hodge classes on self-products of a variety with an automorphism}, Compositio Math.114 (1998), 329--336. Unnumbered main theorem, p.329; polarization correction, p.330; \S7, rational Hermitian classification; Proposition10; \S\S11--13, sixfold Prym dominance. \url{https://doi.org/10.1023/A:1000566205021}.
\bibitem{PetersSterk} C.~Peters and H.~Sterk, \emph{Symmetric and Quadratic Forms}, Corollary3.3.4 (Meyer's theorem). Author-hosted text: \url{https://cpeters1.win.tue.nl/Books/QuadraticForms/QuadForms.pdf}.
\bibitem{Markman2025} E.~Markman, \emph{Cycles on abelian $2n$-folds of Weil type from secant sheaves on abelian $n$-folds}, arXiv:2502.03415, supplied source \texttt{250531.tex}. Theorem1.5.1, source label \texttt{thm-algebraicity}; Corollary1.6.1, under ``Verification of the Hodge conjecture for abelian fourfolds.'' The source statement and reduction are read; its complete proof is a cited input, not independently rebuilt here.
\bibitem{Alpoge2026} L.~Alp\"oge, \emph{Remark}, 20 November2026, supplied four-page note and Lean files. The RH/Dirichlet GRH conclusions are used as instructed. Exact formal statement locators: \texttt{ALPOGE\_Main.lean}, declarations \texttt{riemannHypothesis}, \texttt{LFunction\_ne\_zero\_of\_half\_lt\_re}, \texttt{LFunction\_eq\_zero\_re}, and \texttt{LFunction\_eq\_zero\_re\_of\_isPrimitive}, with definitions in \texttt{ALPOGE\_Statements.lean}. The note's steps2--3 locate its inert-prime/Eisenstein setting. No Lean rebuild is claimed.
\bibitem{Lee2023} E.~S.~Lee, \emph{The prime number theorem for primes in arithmetic progressions at large values}, arXiv:2301.13457v2, Corollary1.1, equation(2), the following unnumbered theta inequality, and Table1, printed p.2. \url{https://arxiv.org/abs/2301.13457}.
\bibitem{Masullo2026} Masullo, \emph{The Music of the Primes}, draft20 November2026, supplied statement index only. Theorem13.3.1.2, ``Divisor of $L(s,\chi)$,'' PDF p.522, is indexed with unfinished statement/proof instructions. General Hecke GRH consequences remain explicitly conditional (Masullo); no such consequence is needed for the displayed geometric bounds.
\bibitem{GrothArchive14} A.~Grothendieck, Montpellier archive, cote14, PDF scans9,38--39,44--47, corresponding to archive folios8,37--38,43--46. The last group is the letter to Serre,27 August1965. \url{https://grothendieck.umontpellier.fr/14.pdf}.
\bibitem{GrothArchive16} A.~Grothendieck, Montpellier archive, cote16: PDF scans44--47, archive folios43--46, letter to Coates dated4 January1967; PDF scans48--50, archive folios47--49, a separate letter to Coates whose handwritten date reads6 January1966. The latter date is reported as written, without emendation. \url{https://grothendieck.umontpellier.fr/16.pdf}.
\end{thebibliography}
